% Successor integration 2026-09-26: the longitudinal reader precedes G.
% The earlier circular realization is retained as a corollary, not a generator.
\section{Action, structural length, gravitation, and thermal information}
\label{v:ant:sec:gravity}

\subsection{Registered return and construction of length}

Write $\alpha_{\HMT}$ for the already generated dodecaphasic coordinate
---denoted $\alpha$ in the preceding sections---. No metrological value of
$G$ or tabulated Planck length enters here. The marked record comprises
\begin{equation}
 N=12\cdot4\cdot120=5760
 \label{v:eq:cardinal-longitud-alpha}
\end{equation}
positions. Let $u_\Omega=N^{-1/2}\mathbf1_\Omega$,
$Q_\Omega=I-u_\Omega u_\Omega^*$, and let
$\mathsf B:V_U\to V_\Omega$ be the typed Hadamard return, with
\begin{equation}
 \mathsf B^*\mathsf B=\mathsf B\mathsf B^*=\frac14I,
 \qquad \mathsf C=Q_\Omega\mathsf B=\mathsf BQ_U.
 \label{v:eq:retorno-hadamard-longitud}
\end{equation}
For the diagonal pinching
$\Delta_\Omega(A)=\sum_{i\in\Omega}E_iAE_i$, one obtains
\begin{equation}
 \Delta_\Omega(\mathsf C\mathsf C^*)
 =\frac{N-1}{4N}I=r_NI,
 \qquad
 r_N=\frac{5759}{23040},
 \label{v:eq:residuo-longitud-alpha}
\end{equation}
whereas the complementary channel retains $1/(4N)$. Indeed,
$\mathsf C\mathsf C^*=Q_\Omega/4$, and every diagonal entry of
$Q_\Omega$ equals $1-1/N$. Diagonal pinching determines the longitudinal
response; the enlarged record and its complement retain the memory that the
diagonal readout does not distinguish.

Let $\mathcal L_L$ be the oriented length line and let $L_*>0$ be its
reference section transported by the longitudinal cochain. This section is
not silently identified with the native edge basis: its relation to the
clock is fixed later by the circular realization. The already constructed
$H_4$-cokernel norm supplies
$\mathrm N_{H_4}(\alpha_{\HMT})=\alpha_{\HMT}^{16}$. Hence the linear
composition of the registered return determines
\begin{equation}
 \boxed{
 L_\alpha=\alpha_{\HMT}^{16}r_NL_*
 =\alpha_{\HMT}^{16}\frac{5759}{23040}L_* .}
 \label{v:eq:longitud-alpha}
\end{equation}
The power, the incidence coefficient, and the dimensional section have
separate and explicit origins. $L_\alpha$ is a longitudinal output prior
to the gravitational readout.

\subsection{Operator composition and unique determination of gravitation}

Let $\mathcal H\ne\{0\}$ be a finite Hilbert fiber and let $Y=UY_0U^*>0$
be the positive invertible character of a route, with its transport, sheet,
orientation, and memory retained. For an action section
$\hbar_a>0$, $a\in\{\mathrm{pre},\mathrm{ret}\}$, and the speed
$c=c_{\mathrm{int}}=c_{\mathrm{vac}}=\widehat c\,u_C>0$ already constructed in
\eqref{v:eq:identidad-velocidad}, set
\begin{equation}
 Q_{L,a}=\frac{\hbar_ac}{L_\alpha},\qquad
 H_a=Q_{L,a}Y,\qquad M_a=\frac{H_a}{c^2},\qquad
 \Lambda_a=L_\alpha Y^{-1},\qquad R_a=L_\alpha Y.
 \label{v:eq:lectores-longitud-accion}
\end{equation}
The null sector is kept separate and is not inverted.

\begin{theorem}[Uniqueness of the gravitational coefficient in the radial realization]
\label{v:thm:gravedad-unica-alpha}
The preceding readers satisfy
\begin{equation}
 H_a\Lambda_a=\hbar_ac\,I,\qquad
 R_a\Lambda_a=L_\alpha^2I.
 \label{v:eq:dos-relaciones-gravedad}
\end{equation}
If the declared physical reading of the reduced radius is
$R_a=(G_a/c^2)M_a$, then there is a unique compatible coefficient, namely
\begin{equation}
 \boxed{G_a=\frac{c^3L_\alpha^2}{\hbar_a}},
 \qquad
 \boxed{\frac{G_a\hbar_a}{c^3}=L_\alpha^2}.
 \label{v:eq:gravedad-unica-alpha}
\end{equation}
\end{theorem}
\begin{proof}
The definitions give
$H_a\Lambda_a=Q_{L,a}L_\alpha I=\hbar_acI$ and
$R_a\Lambda_a=L_\alpha^2I$. Eliminating $\Lambda_a$ yields
\[
 R_a=\frac{L_\alpha^2}{\hbar_ac}H_a.
\]
On the other hand,
$(G_a/c^2)M_a=G_aH_a/c^4$. Since $H_a$ is invertible on the nonzero
fiber, the two expressions agree if and only if
$G_a=c^3L_\alpha^2/\hbar_a$. The same cancellation shows that every
other $G_a'$ preserving both relations satisfies $G_a'=G_a$.
\end{proof}

The proof uses the two readers jointly on the same character; it does not
identify a compression of $Y^{-1}$ with the inverse of a compression. A
route variation or memory elimination acting identically on $H_a$ and
$R_a$ preserves their proportionality and introduces no new scalar factor.
Dimensional analysis confirms the exponents
$c^3L_\alpha^2\hbar_a^{-1}$, but it is
\eqref{v:eq:dos-relaciones-gravedad}, rather than that dimensional check,
that fixes the dimensionless coefficient to one.

\subsection{The circular expression as a corollary}

The calendar determines $T_{108}=108t_0$ and
$\omega_{108}=2\pi_{\HMT}/T_{108}$. In the circular realization of the
same longitudinal reader,
\begin{equation}
 \omega_{108}=\frac{c}{L_\alpha},\qquad
 R_{108}=\frac{c}{\omega_{108}}
 =\frac{54}{\pi_{\HMT}}ct_0=L_\alpha,
 \qquad
 t_0=\frac{\pi_{\HMT}L_\alpha}{54c}.
 \label{v:eq:compatibilidad-reloj-longitud}
\end{equation}
This equality declares compatibility between the lifted clock advance and
the longitudinal section; it does not use $G_a$ to select $L_\alpha$.
Substitution into \eqref{v:eq:gravedad-unica-alpha} gives
\begin{equation}
 G_a=\left(\frac{54}{\pi_{\HMT}}\right)^2
 \frac{c^5t_0^2}{\hbar_a}.
 \label{v:eq:gravedad-circular-corolario}
\end{equation}
Consequently, in the parametrized expression
$G_a(\vartheta)=\vartheta c^5t_0^2/\hbar_a$, the value
$\vartheta=(54/\pi_{\HMT})^2$ is a corollary of the constructed length,
not an independent circular selector.

With
$E_{108,a}=E_{P,a}=Q_{L,a}=\hbar_a\omega_{108}$ and
$m_{108,a}=E_{108,a}/c^2$, one also retains
\[
 \frac{\hbar_a}{m_{108,a}c}=R_{108}=L_\alpha,
 \qquad
 \frac{h_a}{m_{108,a}c}=2\pi_{\HMT}L_\alpha.
\]
Finally, the two action orientations satisfy
\begin{equation}
 \frac{G_{\mathrm{pre}}}{G_{\mathrm{ret}}}
 =\frac{\hbar_{\mathrm{ret}}}{\hbar_{\mathrm{pre}}}
 =R_{\mathrm{act}}^{-1},
 \qquad
 \ell_P^2:=L_\alpha^2
 =\frac{G_a\hbar_a}{c^3}.
 \label{v:eq:reciprocidad-gravedad-accion}
\end{equation}
The areal unit thus originates in $L_\alpha^2$; the later reciprocity
between action and gravitation preserves that area under a change of
orientation.
